\chapter{Background}\label{chap:background}

% --------------------- Section 1 ---------------------

\section{The Rust Programming Language.}\label{sec:rust}
Rust is a system programming language that focuses on safety, speed and concurrency.
It ensures memory safety without garbage collection, meaning pure Rust programs are free from null pointer dereferences and data races.

\paragraph{Ownership and Borrowing}
Rust's ownership system, integrated into its type system, is inspired by \textit{linear logic}~\cite{Girard87, Girard95} and \textit{linear types}~\cite{Wadler90, Odersky92}, enforcing that each piece of memory (a value) has a single \textit{owner} (the variable binding) at any time~\cite{Clarke98, Boyapati02}.
When the owner goes out of scope, the memory is automatically deallocated, enabling user-defined destructors and supporting the \textit{resource acquisition is initialization} (RAII) pattern~\cite{Stroustrup94}.
Ownership can be transferred (\textit{moved}) or temporarily shared (\textit{borrowed}) through references.
Rust supports two types of borrows: multiple \textit{immutable} borrows or a single \textit{mutable} borrow, but never both at the same time.
These constraints, enforced by the compiler through \textit{borrow checking}, guarantee memory safety and prevent dangling pointers by ensuring that reference lifetimes never outlive their owners.
To support low-level operations, Rust provides \inlinerust{unsafe} blocks, where the compiler's safety guarantees are suspended and the burden of avoiding undefined behavior falls on the programmer.
Outside of these blocks, Rust enforces strict safety, making undefined behavior impossible in safe code.

\paragraph{Structs, Traits and Impls}
Structs (\inlinerust{struct}) and enums (\inlinerust{enum}) are the primary composite data types in Rust, enabling the creation of complex data structures with named fields and variant types, respectively.

\begin{figure}[htbp!]
	\centering
	\begin{minipage}{0.49\textwidth}
		\centering
		\showrust[0.98\linewidth]{struct.rs}
		\captionsetup{type=listing,skip=2pt}
		\caption{Example of a Rust struct}
	\end{minipage}
	\hfill
	\begin{minipage}{0.49\textwidth}
		\centering
		\showrust[0.98\linewidth]{enum.rs}
		\captionsetup{type=listing,skip=2pt}
		\caption{Example of a Rust enum}
	\end{minipage}
	\label{fig:structs-enums}
\end{figure}

Rust's trait system is a central mechanism for expressing polymorphism and shared behavior across types~\cite{RustBook,Matsakis14}.
Traits (\inlinerust{trait}) define a set of method signatures that types can implement, enabling abstraction over operations without committing to a concrete type.
Implementation blocks (\inlinerust{impl}) associate specific types with traits, providing concrete method definitions that satisfy the trait's contract.
Rust also supports \textit{inherent implementations}, where methods are defined directly on a type without requiring a trait, allowing type-specific functionality and constructor-like associated functions.

\begin{figure}[htbp!]
	\centering
	\begin{minipage}{0.49\textwidth}
		\centering
		\showrust[0.98\linewidth]{trait-impl.rs}
		\captionsetup{type=listing,skip=2pt}
		\caption{Example of a Rust trait impl}
	\end{minipage}
	\hfill
	\begin{minipage}{0.49\textwidth}
		\centering
		\showrust[0.98\linewidth]{inherent-impl.rs}
		\captionsetup{type=listing,skip=2pt}
		\caption{Example of a Rust inherent impl}
	\end{minipage}
	\label{fig:traits-impls}
\end{figure}

By combining \textit{traits}, \textit{trait bounds} and \textit{impl blocks}, Rust achieves ad hoc polymorphism~\cite{Cardelli85,Wadler90} and generic programming~\cite{Musser89,Backhouse99} while preserving static guarantees and enabling the compiler to resolve method calls efficiently at compile time.
Trait-based design is tightly integrated with Rust's ownership and borrowing model, ensuring that polymorphic abstractions do not compromise memory safety.

\paragraph{Software Variability}
Rust supports software variability through its attribute system, particularly the \inlinerust{cfg} attribute, which allows conditional compilation based on specified configuration predicates
A feature in Rust is defined by all \textit{terms} annotated with the same \textit{config name} in a \inlinerust{cfg} \textit{outer attribute}.
The inclusion of a feature in a product depends on the evaluation of a \textit{configuration predicate} at compile time, e.g. \inlinerust{#[cfg(any(unix, windows))]}. 
\textit{Cargo}, Rust's package manager, lets developers declare features using \textit{config names} in the \texttt{Cargo.toml} file, combining them with logical operators into \textit{configuration predicates}.
Feature dependencies, i.e., \textit{cross-tree constraints} of \textit{feature models}~\cite{Kang90,Seidl16,Karatas10}, can also be specified in the same file.

\paragraph{Release Channels}
Rust is distributed through three release channels: \textit{stable}, \textit{beta} and \textit{nightly}.
These channels differ in maturity guarantees and access to experimental language features~\cite{RustReleaseChannels}.
The \textit{stable} channel prioritizes reliability and backwards compatibility: only fully tested features are enabled, ensuring that libraries and applications compiled on stable continue to work across future releases.
The \textit{beta} channel serves as a staging ground for upcoming stable releases, allowing developers to test new features before they reach stable.
In contrast, the \textit{nightly} channel is updated daily and exposes experimental compiler internals, unstable language constructs and opt-in features (guarded by \inlinerust{#![feature(...)]} attributes).
These features include advanced compiler hooks, low-level intrinsics and incomplete language extensions such as \textit{specialization} and are available exclusively on nightly because their behavior is still under active design and soundness evaluation.
Although nightly enables powerful experimentation, it provides no stability guarantees: features may change, be redesigned or removed without deprecation.
For this reason, production-grade Rust codebases overwhelmingly target the stable channel, limiting developers from relying on certain expressive but unstable mechanisms.

% --------------------- Section 2 ---------------------

\section{Dispatch Techniques}\label{sec:dispatch}
Modern programming languages employ a variety of dispatch strategies at run-time or compile-time to determine which function or method is executed.
These strategies fundamentally influence a language's expressiveness, performance guarantees and capacity for static reasoning~\cite{Cardelli85,Milner78}.
Rust integrates several dispatch paradigms to reconcile expressive abstraction with predictable performance, adhering to the principle of \textit{zero-cost abstraction}~\cite{Steele90,Stroustrup05}.
Most decisions are resolved statically, enabling the compiler to generate optimized machine code without runtime overhead~\cite{Aho06,Cytron91}.
At the same time, Rust offers controlled forms of dynamic dispatch where late binding is essential, such as trait objects~\cite{Matsakis14}.
Beyond these core mechanisms, techniques like inlining heuristics~\cite{Cavazos05}, function cloning~\cite{Das03} and trait-based overloading~\cite{Stuckey05} enable the compiler to further optimize code based on contextual information~\cite{Chang97,Wegman85}.
Together, Rust's dispatch mechanisms provide a flexible yet robust model that preserves safety and performance.

\paragraph{Monomorphization}
Monomorphization~\cite{Griesemer2020FeatherweightG} is the process by which the compiler generates concrete, type-specific versions of generic functions, methods and data structures.
When a generic function is instantiated with a particular type parameter, the compiler produces a specialized version of the function in the final binary.
This approach mirrors C++ template instantiation~\cite{Vandevoorde02}, enabling aggressive optimizations such as inlining, constant propagation and dead-code elimination~\cite{Aho06,Wegman85,Cytron91}.

\begin{figure}[htbp!]
	\centering
	\begin{minipage}{0.49\textwidth}
		\centering
		\showcpp[0.98\linewidth]{monomorphization-pre.cpp}
	\end{minipage}
	\hfill
	\begin{minipage}{0.49\textwidth}
		\centering
		\showcpp[0.98\linewidth]{monomorphization-post.cpp}
	\end{minipage}
	\captionsetup{type=listing,skip=2pt}
	\caption{Monomorphization process in C++: on the left, a generic function \inlinecpp{add<T>} is defined and instantiated with two different types (\inlinecpp{int} and \inlinecpp{double}); on the right, the compiler generates two specialized versions of the function (\inlinecpp{add_int} and \inlinecpp{add_double}) and uses them accordingly}
	\label{fig:monomorphization}
\end{figure}

Because all type information is known at compile time, monomorphized functions avoid virtual dispatch overhead and allow the compiler to reason precisely about memory layout~\cite{Lattner04}.
The primary trade-off is \textit{code bloat}: large applications with many generic instantiations may experience increased binary size.
Nevertheless, Rust's monomorphization strategy is central to its performance model, making generic abstractions effectively as fast as handwritten, type-specific code.

\paragraph{Overloading}
Overloading allows multiple functions or methods to share the same name but differ in parameter types or counts.
Rust does not support traditional function overloading by parameter types as seen in languages like C++ or Java.
Instead, it employs a unified mechanism based on \textit{trait resolution} and \textit{associated functions}~\cite{RustBook}.
Overloading emerges indirectly: different types may provide different implementations of the same trait methods and users may introduce trait bounds to guide method selection.
This form of ad hoc polymorphism resembles \textit{type classes} in Haskell~\cite{Wadler88,Hall96}, where method resolution is determined by the combination of the type and the trait.
Rust enforces coherence to guarantee that the compiler can unambiguously select an implementation at compile time~\cite{Chlipala11,Garrigue98}, preventing conflicts caused by independently defined overlapping instances.
In addition, traits can define associated functions that behave similarly to static methods, allowing for API designs that mimic overloading while remaining type-safe and deterministic.
This approach ensures predictable dispatch behavior, maintains clarity in method resolution and avoids the complexities of multi-parameter or ambiguous overload sets common in other languages~\cite{PeytonJones00}.

\paragraph{Specialization}
Specialization is an advanced dispatch mechanism that allows more specific trait implementations to refine or override default behavior.
In Rust, specialization aims to resolve long-standing expressiveness gaps in the trait system by enabling prioritized selection among overlapping implementations.
This is particularly useful when a general-purpose implementation exists (often expressed as a \textit{default} method in a trait) but certain types require optimized or semantically distinct behavior~\cite{SpecializationRFC}.
The feature builds on Rust's \textit{coherence} and \textit{orphan rules}\cite{RustOrphanRules}, ensuring that overlapping implementations maintain type-safety and global consistency~\cite{Chlipala11}.
Although as of today specialization remains unstable and experimental (through the \inlinerust{#![feature(specialization)]} or \inlinerust{#![feature(min_specialization)]} flags~\cite{SpecializationRFC,MinSpecialization}), it provides a foundation for advanced library design patterns, such as optimized container types, numerical traits and fallback-based APIs.
The mechanism is conceptually related to multi-methods and ad hoc polymorphism in other languages~\cite{Chambers92,Castagna95}, but Rust restricts its form to preserve determinism and avoid ambiguity in resolution that have yet to be fully addressed.

\begin{figure}[htbp!]
	\centering
	\showcpp*[0.98\linewidth]{specialization.cpp}
	\captionsetup{type=listing,skip=2pt}
	\caption{Specialization in C++. The code is the equivalent of that in Listing~\ref{fig:monomorphic_specialization_pre}}
	\label{fig:cpp-specialization}
\end{figure}

\begin{figure}[htbp!]
	\centering
	\showrust*[0.98\linewidth]{unstable-specialization.rs}
	\captionsetup{type=listing,skip=2pt}
	\caption{Specialization in Rust (unstable). The code is the equivalent of that in Listing~\ref{fig:cpp-specialization}}
	\label{fig:rust-specialization}
\end{figure}

% --------------------- Section 3 ---------------------

\section{Metaprogramming}\label{sec:metaprogramming}
Metaprogramming refers to the practice of writing code that generates or manipulates other code, enabling abstraction, automation and domain-specific optimizations.
Historically rooted in Lisp's symbolic macro systems~\cite{McCarthy60} and later expanded through compile-time techniques such as C++ templates~\cite{Vandevoorde02}, metaprogramming has evolved into a fundamental paradigm in modern language design, as highlighted by many surveys and studies~\cite{MetaprogrammingLanguagesSurvey19, DSLMetaProgramming08}.
Its goals are multifaceted: reducing boilerplate, enforcing structural invariants, encoding repetitive patterns and enabling compile-time specialization or reasoning that would otherwise be impossible or inefficient at run-time.
Across programming ecosystems, metaprogramming serves both practical and expressive purposes, ranging from code generation and DSL construction to advanced type-level computation~\cite{DSLMetaProgramming08, MetaprogrammingUr10}.
Rust approaches metaprogramming through a powerful and expressive macro system designed to maintain the language's guarantees of safety, determinism and predictable compilation.
Rust distinguishes between two primary categories: \textit{declarative macros} and \textit{procedural macros}.

\paragraph{Declarative Macros}
Declarative macros (defined using \inlinerust{macro_rules!}) operate using pattern matching over token streams, allowing developers to define custom syntactic constructs that expand into valid Rust code.
While reminiscent of Lisp macros in spirit, Rust's declarative macros enforce \textit{hygiene}~\cite{HygienicMacros92,Kohlbecker1986HygienicME,Clinger2020HygienicMT}: macro-generated identifiers cannot inadvertently capture variables or introduce unintended bindings.
This prevents entire classes of errors common in textual or non-hygienic macro systems.

\paragraph{Procedural Macros}
Procedural macros provide direct access to the compiler's abstract syntax tree (AST~\cite{Aho06,CloneDetection98}) through the \inlinerust{proc_macro} API.
These macros behave as compiler plugins that transform annotated code items at compile time.
Three major forms exist: \textit{function-like procedural macros}, \textit{attribute macros} and \textit{derive macros}.
The latter are widely used to automatically implement traits such as \inlinerust{Clone}, \inlinerust{Debug} or \inlinerust{Serialize}, significantly reducing repetitive boilerplate and enabling library authors to design ergonomic declarative APIs.
Attribute macros allow developers to attach custom semantics or transformations to items, facilitating the creation of embedded domain-specific languages (DSLs) and compile-time validation tools; the \inlinerust{#[when(...)]} attribute macro introduced in this thesis is an example of such usage.
Function-like procedural macros (e.g. the \inlinerust{spec!} macro) enable arbitrary code generation based on input token streams, allowing for complex compile-time computations and code synthesis.
Because procedural macros operate on structured syntax trees instead of raw text, they can enforce well-formedness and preserve strong static guarantees during expansion.

\paragraph{Integration with the Compilation Pipeline}
Rust's macro system integrates deeply with the compiler's type checking and optimization phases.
Unlike languages where macros expand at a separate preprocessing stage (e.g. the C preprocessor~\cite{Dietrich2021CppSigET}), Rust expands macros as part of the front-end compilation pipeline, ensuring that hygiene, scoping and type resolution remain consistent.
This design enables robust static analyses to occur alongside macro expansion and avoids undesired ambiguities or semantic inconsistencies.
The combination of hygienic rewriting, AST-level transformations and compile-time type awareness makes Rust's metaprogramming capabilities both expressive and safe, empowering developers to construct efficient abstractions without compromising reliability.
